\subsection{A surface counterexample with a smooth very ample section}

\begin{lemma}[A smooth root cover makes a generated polarization very ample]
\label{lem:OC-Kummer-very-ample}
\status{Proved; independently audited by the root and reduction-descent agents.}
Let $S$ be a smooth projective connected surface over an algebraically closed field of characteristic $p$.  Let $L$ be ample and generated by global sections, and choose an integer $q>1$, prime to $p$, such that $L^q$ is very ample.  Then there is a smooth projective connected surface $X$, a finite flat map $\rho:X\to S$ of degree prime to $p$, and a very ample line bundle $A=\rho^*L$.  Moreover $H^1(S,L^{-1})$ is a direct summand of $H^1(X,A^{-1})$.
\end{lemma}
\begin{proof}
Choose a basis $b_1,\ldots,b_r$ of $H^0(S,L^q)$ sufficiently general that its zero divisors $D_i$ are smooth irreducible curves, pairwise transverse, with no triple intersections.  Such bases form a nonempty open subset of the space of bases: apply Bertini successively, avoiding the finitely many previous pairwise intersection points.  Each smooth ample divisor is connected, hence irreducible here.

Form the fibre product over $S$ of the $r$ cyclic covers
\[
 X=\operatorname{Spec}_S\left(
  \bigotimes_{i=1}^r\left(\bigoplus_{j=0}^{q-1}L^{-j}\right)\right),
 \qquad y_i^q=b_i.
\]
Its algebra is locally free of rank $q^r$, and $y_i$ is a global section of $\rho^*L$.  At a point lying on exactly $a$ of the $D_i$, where $0\leq a\leq2$, their local equations are part of a regular system of parameters.  Adjoining their $q$th roots replaces those parameters by roots; adjoining the remaining roots is \'etale because their equations are units and $q$ is invertible.  Thus $X$ is smooth.

We also check connectedness.  Choose a rational trivialization $s$ of $L$ whose divisor contains none of the $D_i$ and set $c_i=b_i/s^q$.  At the discrete valuation of $D_i$, one has $v(c_i)=1$ and $v(c_j)=0$ for $j\ne i$.  In the tower adjoining the roots in order, every earlier extension is unramified at this valuation: it is obtained by adjoining roots of units with $q$ invertible.  At any prolongation of the valuation the polynomial $T^q-c_i$ is Eisenstein.  Hence its degree is $q$ at each stage, including for composite $q$.  The generic algebra is therefore a field of degree $q^r$.  Since $X/S$ is flat and $X$ is smooth, every component dominates $S$; thus $X$ is integral and connected.

Let $s_0,\ldots,s_m$ be a basis generating $L$.  The sections
\[
                 \rho^*s_0,\ldots,\rho^*s_m,y_1,\ldots,y_r
\]
define a morphism to projective space.  On the chart $s_\alpha\ne0$, the affine coordinate ring of $S_{s_\alpha}$ is generated by the functions $b_i/s_\alpha^q$, because $L^q$ is very ample and the $b_i$ form a basis.  The algebra of $X$ on this chart is generated over it by $y_i/s_\alpha$, and
\[
             b_i/s_\alpha^q=(y_i/s_\alpha)^q.
\]
Consequently the displayed projective morphism is a closed immersion on each of these charts.  They cover $X$ because the $s_\alpha$ generate $L$.  Properness then makes the morphism a closed immersion globally.  Its hyperplane bundle is $A=\rho^*L$, proving very ampleness.

Finally $\rho_*\OO_X$ is a direct sum of negative powers of $L$, with one summand $\OO_S$.  Projection formula and the finiteness of $\rho$ put $H^1(S,L^{-1})$ as a direct summand in $H^1(X,A^{-1})$.
\end{proof}

\begin{lemma}[An explicit Tango curve with a generated root]
\label{lem:OC-generated-Tango-root}
\status{Proved under stated hypotheses}
For every prime $p$, put
\[
 h=\begin{cases}8,&p=2,\\p^2-2,&p\geq3,\end{cases}
 \qquad m=hp-1,\qquad d=(p-1)h-2,\qquad \ell=p+1.
\]
Let $C$ be the smooth projective model of
\[
                        y^p-y=x^m.
\]
There is a unique point $P_\infty$ above infinity, and
\[
 (dx)=pdP_\infty,\qquad 2g(C)-2=pd.
\]
The positive divisor $D=dP_\infty$ is a Tango divisor.  It has the root
\[
 N=\OO_C\bigl((d/\ell)P_\infty\bigr),\qquad
 N^\ell=\OO_C(D)=:T,
\]
and $N$ is generated by global sections.  For $p=2$ one has
$g=7$ and $N=\OO_C(2P_\infty)$; for odd $p$ one has
$d=p(p+1)(p-2)$ and $N=\OO_C(p(p-2)P_\infty)$.
\end{lemma}
\begin{proof}
The Artin--Schreier polynomial is irreducible over $k(x)$: a rational function $a$ with $a^p-a=x^m$ would have no finite pole and would have a pole at infinity whose order, after applying $a^p-a$, is divisible by $p$, whereas $p\nmid m$.  For an Artin--Schreier polynomial of prime degree, its Galois group is a subgroup of the translation group $\Fp$, so absence of a root makes it irreducible and proves the asserted extension degree.  The affine curve is smooth and \'etale over the affine $x$-line, since the derivative with respect to $y$ is $-1$.

At a place over infinity the relation forces $p\,v(y)=m\,v(x)$, so its ramification index is divisible by $p$.  Thus it is $p$, there is one such place, and
$v(x)=-p$, $v(y)=-m$.  Put
\[
                    w=y/x^h,\qquad v=1/x,\qquad c=h(p-1).
\]
Then $w$ has valuation one, and the local equation is
\[
                         v+w v^c=w^p.
\]
Its derivative with respect to $v$ is a unit at $(w,v)=(0,0)$, so this is the smooth local model with parameter $w$, and $v$ has valuation $p$.  Differentiating gives
\[
             dx=\frac{v^{c-2}}{1+cw v^{c-1}}\,dw.
\]
The denominator is a unit.  Since $dx$ has no zero or pole on the affine curve, its full divisor is $p(c-2)P_\infty=pdP_\infty$.  This also computes the genus and proves that $D$ is a Tango divisor in the sense of \cite[Definition~2.1]{Zheng2017Kodaira}, equivalently \cite[Theorem~1.2]{Mukai2011Kodaira}.

The displayed choices give $d/\ell=2$ at $p=2$, and $d/\ell=p(p-2)$ for odd $p$.  The functions $1,x$ in the first case, and $1,x^{p-2}$ in the second, generate $N$: the first section is nonzero off $P_\infty$ and the second has exactly the pole order prescribed by $N$, so is nonzero there after trivialization.
\end{proof}

\begin{proposition}[A generated ample Raynaud polarization with nonvanishing]
\label{prop:OC-generated-Raynaud}
\status{Proved; independently audited by the root and reduction-descent agents.}
For every prime $p$, the curve in Lemma~\ref{lem:OC-generated-Tango-root} yields an explicitly specified smooth projective connected surface $S$ and an ample globally generated line bundle $L$ such that
\[
                         H^1(S,L^{-1})\ne0.
\]
Moreover $L^{(p+1)^2}$ is very ample.
\end{proposition}
\begin{proof}
Use $T=N^\ell=\OO_C(D)$ from the preceding lemma.  The exact differential $dx$ gives the nonzero class in
\[
 \ker\bigl(H^1(C,T^{-1})\xrightarrow{F}H^1(C,T^{-p})\bigr)
\]
by \cite[Theorem~1.2 and its proof in~1.2]{Mukai2011Kodaira}.  It determines an extension
\[
              0\longrightarrow\OO_C\longrightarrow E
                  \longrightarrow T\longrightarrow0.
\]
The projective bundle $\pi:P=\mathbf P_C(E)\to C$ parametrizes one-dimensional quotients.  Let $E_0$ be the section from $E\twoheadrightarrow T$.  Frobenius splitting of the extension gives the disjoint purely inseparable degree-$p$ multisection $G$.  The precise construction and divisor formulas are
\cite[Section~2.2, Proposition~2.2]{Zheng2017Kodaira}:
\[
 \OO_P(E_0)=\OO_P(1),\qquad
 \OO_P(G)=\OO_P(p)\otimes\pi^*T^{-p},\qquad E_0\cap G=\varnothing.
\]
Both curves are smooth.  To verify the smoothness input, locally the inseparable multisection has equation $z^p=f_i$, where $df_i$ is the local expression of $dx$ after dividing by a local generator of $T^p$.  The divisor identity $(dx)=pD$ says precisely that $df_i$ is nowhere zero.  The Jacobian criterion proves smoothness, as in the complete proof of \cite[Proposition~1.7]{Mukai2011Kodaira}.  Thus no assertion about the smoothness of a purely inseparable cover is assumed without this check.

Put
\[
 M=\OO_P(-1)\otimes\pi^*N^p,
 \qquad M^{-\ell}=\OO_P(E_0+G),
\]
and take the degree-$\ell$ cyclic cover
\[
 \psi:S=\operatorname{Spec}_P\!\left(\bigoplus_{i=0}^{\ell-1}M^i\right)
       \longrightarrow P
\]
branched on $E_0+G$.  The branch curves are disjoint and smooth, and $\ell=p+1$ is prime to $p$.  The local cyclic equations therefore prove that $S$ is smooth.  The valuation at $E_0$ makes its generic polynomial Eisenstein, so $S$ is connected and integral.  Let $\widetilde E$ be the reduced inverse image of $E_0$ and set $f=\pi\psi$.  Then
\[
 \psi^*E_0=\ell\widetilde E,
 \qquad \widetilde E^2=\deg N>0,
 \qquad L=\OO_S(p\widetilde E)\otimes f^*N.
\]
Every geometric fibre of $f$ is integral: on its affine chart the equation is $u^\ell=(z-a)^p$, which is irreducible since $\gcd(\ell,p)=1$, and its point at infinity is the smooth section $\widetilde E$.  In particular $\widetilde E$ has intersection one with a fibre.  The effective section $\widetilde E$ is nef, since its self-intersection is positive and its intersection with any other irreducible curve is nonnegative.  On every vertical curve $L$ has positive degree, and on every horizontal curve the summand $f^*N$ has positive degree.  Also $L^2=p(p+2)\deg N>0$.  Nakai--Moishezon proves that $L$ is ample.

We now prove global generation, the extra fact required beyond the original very-ampleness obstruction.  The cyclic eigensheaf decomposition is
\begin{equation}\label{eq:OC-positive-eigensheaves}
 \psi_*\OO_S(p\widetilde E)
       =\OO_P\oplus\bigoplus_{i=1}^p M^i(E_0).
\end{equation}
Indeed, $p=\ell-1$ and
$\OO_S(p\widetilde E)=\OO_S(-\widetilde E)\otimes\psi^*\OO_P(E_0)$;
for the negative divisor the index-zero eigensheaf is $\OO_P(-E_0)$ and all other eigensheaves are $M^i$.  This is also the $k=1$ case of
\cite[Proposition~3.3]{Zheng2017Kodaira}.  Since
$M^i(E_0)=\OO_P(1-i)\otimes\pi^*N^{pi}$, pushing \eqref{eq:OC-positive-eigensheaves} to $C$ gives
\[
                       f_*L=N\oplus N^{p+1}.
\]
The index-zero summand generates $L$ away from $\widetilde E$, because it is the canonical section of $\OO_S(p\widetilde E)$ multiplied by sections generating $N$.  On $\widetilde E$, the index-one summand generates: if $s=0$ is a local equation of $E_0$ and $u$ is the cyclic parameter, then $u^\ell=s$ up to a unit, and the generator $u/s$ of $M(E_0)$ is a unit times $u^{-p}$.  It is consequently a local generator of $\OO_S(p\widetilde E)$ along $\widetilde E$.  The coefficient bundle is $N^{p+1}$, which is generated because $N$ is.  This proves global generation everywhere.

For nonvanishing, the negative decomposition gives
\[
 \psi_*L^{-1}
   =\pi^*N^{-1}\otimes
      \left(\bigoplus_{i=0}^{p-1}M^i(-E_0)\ \oplus\ M^p\right).
\]
Every summand has negative degree on a $\pi$-fibre, so $f_*L^{-1}=0$.  Relative duality for the projective line bundle gives
\[
 R^1\pi_*\OO_P(-r)=\bigl(\operatorname{Sym}^{r-2}E\bigr)^\vee\otimes T^{-1}
       \quad(r\geq2).
\]
Here the dual is taken after forming the symmetric power; it is not identified with the symmetric power of the dual in characteristic $p$.  In particular the last eigensheaf contributes the direct summand
\[
 \bigl(\operatorname{Sym}^{p-2}E\bigr)^\vee
          \otimes N^{p^2-p-2}
    =\bigl(\operatorname{Sym}^{p-2}E\bigr)^\vee\otimes T^{p-2}
\]
to $R^1f_*L^{-1}$.  The quotient $E\twoheadrightarrow T$ induces a surjection of symmetric powers, whose dual puts $\OO_C$ into this summand.  It therefore has a nonzero global section.  Leray over the curve yields
\[
        H^1(S,L^{-1})=H^0(C,R^1f_*L^{-1})\ne0.
\]
This also verifies directly the relevant endpoint $a=p$, $b=1$ of
\cite[Theorem~3.7]{Zheng2017Kodaira}; no uncorrected assertion from the earlier Takayama paper is used.

It remains to give a definite very ample power.  Let $q=\ell^2=(p+1)^2$.  Then
\[
 L^q=\psi^*F,\qquad
 F=\OO_P(p\ell)\otimes\pi^*N^{\ell^2}.
\]
The bundle $\operatorname{Sym}^aE\otimes N^b$ on $C$ has a filtration with quotients $N^{b+\ell j}$ for $0\leq j\leq a$.  Whenever $b\deg N\geq2g+1$, all these quotients have vanishing $H^1$ after subtracting any length-two subscheme of $C$, by Serre duality.  The same holds for the filtered bundle.  Its global sections therefore supply arbitrary values and first jets on the base; on a fibre, $\OO_P(a)$ supplies the complete Veronese system.  For $a\geq1$ these two facts separate all length-two subschemes of $P$, proving that $\OO_P(a)\otimes\pi^*N^b$ is very ample.

In the present situation $2g-2=p\ell\deg N$, whereas
\[
 \ell^2\deg N-(2g+1)=\ell\deg N-3>0.
\]
Hence both $F$ and $F\otimes M=\OO_P(p\ell-1)\otimes\pi^*N^{\ell^2+p}$ are very ample, in particular globally generated.  For a cyclic cover, this implies $\psi^*F$ is very ample: sections of $F$ recover the embedded base, and the index-one sections $H^0(P,F\otimes M)$ recover the cyclic coordinate on a neighbourhood of every point, since $F\otimes M$ is generated.  On affine charts this is a surjection onto the cyclic algebra, and properness gives a closed immersion.  Thus the explicit power $L^q$ is very ample, as asserted.
\end{proof}

\begin{theorem}[Full coherent hyperplane restriction fails on smooth surfaces in every characteristic]
\label{thm:OC-surface-coherent-counterexample}
\status{Proved; independently audited by the root and reduction-descent agents.}
For every prime $p$, there is a smooth projective connected surface
$X\subset\mathbf P^r_k$ over an algebraically closed field $k$ of characteristic $p$ and a smooth hyperplane section $Y$ such that
\[
             \ker\bigl(H^1(X,\OO_X)\longrightarrow H^1(Y,\OO_Y)\bigr)\ne0.
\]
The construction is given explicitly by the Artin--Schreier curve, Tango extension, cyclic surface, and finite root-cover equations in the preceding results, with root exponent $q=(p+1)^2$.
\end{theorem}
\begin{proof}
Apply Lemma~\ref{lem:OC-Kummer-very-ample} to the pair $(S,L)$ of
Proposition~\ref{prop:OC-generated-Raynaud}, with $q=(p+1)^2$.  It gives a smooth connected surface $X$ and a very ample $A$ with
$H^1(X,A^{-1})\ne0$.
Choose a general smooth hyperplane section $Y\in|A|$.  It is connected because it is an ample divisor on the smooth projective surface.  Consequently
$H^0(X,\OO_X)=k=H^0(Y,\OO_Y)$, and the exact sequence
\[
       0\longrightarrow A^{-1}\longrightarrow\OO_X
                       \longrightarrow\OO_Y\longrightarrow0
\]
identifies the restriction kernel with $H^1(X,A^{-1})$, which is nonzero.  This proves an actual restriction kernel, rather than negative-line-bundle nonvanishing alone.  In fact the same conclusion holds for every smooth member of $|A|$.

For a fully specified generic choice, first carry out the Artin--Schreier, Tango, and degree-$p+1$ constructions over $\overline{\Fp}$.  Choose any basis of $H^0(S,L^q)$ and apply an invertible matrix of independent indeterminates to it.  Work over the algebraic closure of that rational function field.  Changing the initial basis only makes an invertible linear change in the independent generic matrix coordinates, so this generic construction is unambiguous up to that change of coefficient field coordinates.  The resulting basis lies in the nonempty open set of transverse bases used in the root-cover lemma.  A further generic linear combination of the displayed embedding sections specifies $Y$.  Thus all covers, their exponents, and the required open conditions are fixed by finite equations and generic coefficients.
\end{proof}

\begin{corollary}[The exact coherent degree threshold in characteristic two]
\label{cor:OC-surface-exact-degrees}
\status{Proved; independently audited by the root and reduction-descent agents.}
In characteristic $2$, the preceding construction can be taken with
\[
 C:\ y^2+y=x^{15},\qquad g(C)=7,\qquad
 N=\OO_C(2P_\infty),\qquad \ell=3,\qquad q=9.
\]
For its Raynaud polarization $L$ one has
\[
 h^1(S,L^{-1})=1,
 \qquad H^1(S,L^{-r})=0\quad(r\geq2),
 \qquad h^0(S,L)=6,
 \qquad h^0(S,L^9)=556.
\]
The root cover therefore has degree $9^{556}$ and the displayed sections embed $X$ in $\mathbf P^{561}$.  For every smooth degree-$d$ hypersurface section $Y_d$ in this embedding,
\[
 \dim_k\ker\bigl(H^1(X,\OO_X)\to H^1(Y_d,\OO_{Y_d})\bigr)
   =\begin{cases}1,&d=1,\\0,&d\geq2.\end{cases}
\]
Thus this example has an exact degree-one failure and injectivity in every higher degree.
\end{corollary}
\begin{proof}
For $r\geq1$ and $j=0,1,2$, put
\[
                     m_j=j+\left\lceil\frac{2r-j}{3}\right\rceil.
\]
The $j$th eigensheaf of $\psi_*L^{-r}$ is
$\OO_P(-m_j)\otimes\pi^*N^{2j-r}$.  Each $m_j$ is positive, so $f_*L^{-r}=0$.  The corresponding nonzero higher direct images, with terms of negative symmetric degree omitted, are
\[
 \bigl(\operatorname{Sym}^{m_j-2}E\bigr)^\vee\otimes N^{2j-r-3}.
\]
For $r\geq2$, the largest exponent $2j-r-3$ is $1-r<0$.  The bundle $E^\vee$ has a filtration with quotients $N^{-3}$ and $\OO_C$, so every line quotient of each displayed symmetric-power twist has negative degree.  All its global sections vanish, and Leray gives $H^1(S,L^{-r})=0$.
For $r=1$, the three contributions are, respectively, zero, $N^{-2}$, and $\OO_C$.  Thus $h^1(S,L^{-1})=1$ exactly.

The pole semigroup of $C\setminus\{P_\infty\}$ is generated by $2$ and $15$: the affine algebra is free over $k[x]$ on $1,y$, whose monomials have distinct pole orders $2a$ and $2a+15$.  Hence $h^0(N)=2$, $h^0(N^3)=4$, and the earlier formula $f_*L=N\oplus N^3$ gives $h^0(L)=6$.
The three eigensheaves for $L^9$ push to
\[
 \operatorname{Sym}^{6}E\otimes N^9,\qquad
 \operatorname{Sym}^{5}E\otimes N^{11},\qquad
 \operatorname{Sym}^{4}E\otimes N^{13}.
\]
Their filtrations have line quotients $N^{9+3a}$, $N^{11+3a}$, and $N^{13+3a}$, respectively.  All degrees exceed $2g-2=12$, so Riemann--Roch and additivity of global sections give
\[
 h^0(L^9)
 =\sum_{a=0}^{6}(12+6a)
   +\sum_{a=0}^{5}(16+6a)
   +\sum_{a=0}^{4}(20+6a)
 =210+186+160=556.
\]

For the root cover, projection formula gives
\[
 H^1(X,A^{-d})
  =\bigoplus_{0\leq a_i<9}
          H^1\left(S,L^{-(d+\sum_{i=1}^{556}a_i)}\right).
\]
Only $d=1$ and the unique all-zero tuple contribute, giving dimension one; for $d\geq2$ the group is zero.  Every smooth hypersurface section $Y_d\in|A^d|$ is connected, so the same exact sequence as in Theorem~\ref{thm:OC-surface-coherent-counterexample} identifies this group with the restriction kernel.  The $6+556$ sections from the root-cover lemma give the asserted projective embedding.
\end{proof}
